\chapter{Nonrenormalization and Holomorphy}
\label{ch:susy-nonrenormalization-holomorphy}

\ConventionsLorentzian

Nonrenormalization theorems are statements about Wilsonian coordinates in a
scheme with declared supersymmetry, gauge symmetry, locality, and anomaly
data.  They are not statements about arbitrary effective actions or about
singular infrared limits.  This chapter formulates the perturbative
superpotential theorem and the holomorphic gauge-coupling statement in the
language of Chapter~\ref{ch:susy-wilsonian-schemes-exact-dynamics}.

\section{Chiral Local Coordinates}

Let \(\Phi^i\) be chiral superfields and \(t^A\) chiral coupling coordinates.
A Wilsonian chiral coordinate is a local functional of the form
\[
  I_F[\Phi;t]
  =
  \int d^4x\,d^2\theta\,
  \mathcal W(\Phi;t)+\text{c.c.}
\]
defined modulo \(\bar D_{\dot\alpha}\)-exact full-superspace terms and
equations of motion appropriate to the chosen Wilsonian action.  A full
superspace coordinate is
\[
  I_D[\Phi,\bar\Phi;t,\bar t]
  =
  \int d^4x\,d^4\theta\,
  \mathcal K(\Phi,\bar\Phi;t,\bar t).
\]
The distinction is cohomological: an \(F\)-term is represented by a chiral
integrand and the chiral measure, while a \(D\)-term is represented by the
full superspace measure.

\section{Perturbative Superpotential Theorem}

\begin{theorem}[Perturbative Wilsonian superpotential theorem]
\label{thm:perturbative-wilsonian-superpotential}
Assume a massive or Wilsonian infrared-regulated \(\mathcal N=1\) theory in a
holomorphic Wilsonian scheme.  Perturbative loop corrections do not generate
new superpotential terms beyond those already present as holomorphic local
coordinates, except through tree-level elimination of massive chiral fields
and through explicitly declared anomaly-induced coordinate changes.
\end{theorem}

\begin{proof}
In supergraph perturbation theory each internal propagator carries spinor
derivatives and each interaction vertex is integrated over either \(d^4\theta\)
or over a chiral measure that may be represented inside full superspace by
spinor derivatives acting on delta functions.  For a loop correction to
produce an \(F\)-term, all but one pair of Grassmann integrations would have
to be saturated while leaving only \(d^2\theta\) for the external local
functional.  The spinor-derivative algebra in a loop supplies the full
superspace measure unless an infrared singularity converts a nonlocal
full-superspace expression into a local chiral expression.  A Wilsonian
effective action at scale \(\Lambda\) excludes such singularities by
integrating out only modes above \(\Lambda\).  Therefore loop corrections are
full-superspace local coordinates.  Tree-level elimination of massive chiral
fields is different: it substitutes the solution of a chiral equation of
motion into the existing superpotential and hence stays within the chiral
coordinate class.  Anomaly-induced coordinate changes are precisely the
declared exceptions in the scheme data.
\end{proof}

\section{Spurion Symmetries}

Treat a coupling \(t^A\) as a background chiral superfield.  If an exact
symmetry acts by
\[
  \Phi^i\mapsto e^{\ii q^i\alpha}\Phi^i,\qquad
  t^A\mapsto e^{\ii q^A\alpha}t^A,
\]
then a Wilsonian \(F\)-term must be a holomorphic section with the same total
charge as the chiral measure.  This turns dimensional analysis and
symmetries into a finite problem: list holomorphic monomials in fields and
couplings with the required charges and locality degree.  The conclusion is
only as strong as the declared symmetry and anomaly ledger.

\section{Holomorphic Gauge Coupling}

For a gauge theory with chiral matter, the holomorphic gauge coordinate is
\[
  \frac{1}{16\pi\ii}\int d^4x\,d^2\theta\,
  \tau(\Lambda)\operatorname{tr}(W^\alpha W_\alpha)+\text{c.c.}
\]
In a holomorphic Wilsonian scheme, perturbative corrections to
\(\tau\) are one-loop exact:
\[
  \tau(\Lambda)
  =
  \tau(\Lambda_0)
  +\frac{b_0}{2\pi\ii}\log\frac{\Lambda}{\Lambda_0},
\]
where, in the trace convention of Chapter~\ref{ch:susy-gauge-theory},
\(b_0\) is computed from the adjoint and matter representation indices in
that same bilinear form.

\begin{proof}[Reason for one-loop exactness]
The coordinate \(\tau\) is chiral and dimensionless.  Higher-loop
perturbative corrections would be holomorphic functions of \(\tau\) and the
chiral matter couplings with the same anomalous symmetry transformation as
\(\tau\).  The anomalous rescaling of the charged fields fixes the logarithmic
one-loop shift.  Additional perturbative terms would either violate the
spurion transformation law or be removable by a holomorphic redefinition of
the Wilsonian coordinate.  Physical gauge couplings obtained after canonical
normalization are different coordinates; their beta functions include the
Jacobian of field rescalings and therefore need not be one-loop exact.
\end{proof}

\section{Konishi-Type Rescaling Anomaly}

A chiral field redefinition
\[
  \Phi\mapsto e^{\alpha}\Phi
\]
changes the functional measure by a local chiral anomaly.  In a regulated
scheme this is the statement that the Jacobian contributes
\[
  \Delta S
  =
  \frac{T(R)\alpha}{8\pi^2}
  \int d^4x\,d^2\theta\,\operatorname{tr}(W^\alpha W_\alpha)
  +\text{c.c.}
\]
with \(T(R)\) normalized by the same trace convention as the gauge kinetic
term.  This anomaly explains why holomorphic Wilsonian couplings and
canonically normalized couplings obey different coordinate beta functions.

\section{Status Boundary}

The perturbative theorem above is a Wilsonian theorem under explicit
regularity hypotheses.  Nonperturbative superpotentials, instanton effects,
gaugino condensation, and strong-coupling moduli-space constraints are not
contradictions: they are additional semiclassical or nonperturbative
contributions to Wilsonian chiral coordinates.  Each such contribution must
be derived from a declared integration cycle, zero-mode measure, anomaly
matching, and infrared effective description.

\begin{openproblem}[Rigorous supersymmetric Wilsonian holomorphy]
\label{op:rigorous-susy-holomorphy}
Construct a regulator and BV Wilsonian scheme in which the chiral local
cohomology, supergraph measure, rescaling anomaly, and holomorphic gauge
coordinate are all defined at finite cutoff, and prove the perturbative
nonrenormalization statements as theorems in that scheme.
\end{openproblem}
